\documentclass[12pt,a4paper]{report}

% =====================================================
% GÓI HỖ TRỢ TIẾNG VIỆT
% =====================================================
\usepackage[utf8]{inputenc}
\usepackage[T5]{fontenc}
\usepackage[vietnamese]{babel}

% =====================================================
% ĐỊNH DẠNG TRANG
% =====================================================
\usepackage[a4paper,top=2.5cm,bottom=2.5cm,left=3cm,right=2cm]{geometry}
\usepackage{setspace}
\onehalfspacing
\usepackage{indentfirst}
\setlength{\parindent}{1.5em}

% =====================================================
% ĐỒ HỌA, BẢNG BIỂU
% =====================================================
\usepackage[table]{xcolor}
\usepackage{graphicx}
\usepackage{array}
\usepackage{booktabs}
\usepackage{multirow}
\usepackage{caption}
\usepackage{float}
\usepackage{tabularx}

% =====================================================
% TOÁN HỌC
% =====================================================
\usepackage{amsmath}

% =====================================================
% MỤC LỤC, LIÊN KẾT
% =====================================================
\usepackage[titles]{tocloft}
\usepackage[
    colorlinks=true,
    linkcolor=black,
    urlcolor=blue,
    citecolor=blue
]{hyperref}

% =====================================================
% ĐỊNH DẠNG CHƯƠNG
% =====================================================
\usepackage{titlesec}

% Định dạng số thứ tự chương (La Mã) và phần (Ả Rập)
\renewcommand{\thechapter}{\Roman{chapter}}
\renewcommand{\thesection}{\arabic{chapter}.\arabic{section}}
\renewcommand{\thesubsection}{\thesection.\arabic{subsection}}

\titleformat{\chapter}[block]
{\normalfont\bfseries\large\centering}
{\MakeUppercase{\chaptertitlename}\ \thechapter.}
{0.5em}
{\MakeUppercase}

\titleformat{\section}[block]
{\normalfont\bfseries\normalsize}
{\thesection.\ }
{0.5em}
{}

\titleformat{\subsection}[block]
{\normalfont\bfseries\normalsize}
{\thesubsection.\ }
{0.5em}
{}

% =====================================================
% MỤC LỤC (CẤU HÌNH TOCLOFT)
% =====================================================
\renewcommand{\cftchapfont}{\bfseries}
\renewcommand{\cftchappagefont}{\bfseries}
\renewcommand{\cftchapleader}{\cftdotfill{\cftdotsep}} % Hiện dấu chấm lửng cho chương

% Định dạng chương trong mục lục: Chương I. TÊN CHƯƠNG
\renewcommand{\cftchappresnum}{Chương\ }
\renewcommand{\cftchapaftersnum}{.}
\setlength{\cftchapnumwidth}{2.5cm}

% Định dạng phần trong mục lục: 1.1. Tên phần
\renewcommand{\cftsecaftersnum}{.}
\setlength{\cftsecnumwidth}{1.2cm}

\renewcommand{\cftsubsecaftersnum}{.}
\setlength{\cftsubsecnumwidth}{1.5cm}

\setlength{\cftfignumwidth}{1.8cm}
\setlength{\cfttabnumwidth}{1.8cm}

% Lệnh tùy chỉnh để gom hình và bảng chung một trang
\makeatletter
\newcommand{\customlistoffigures}{%
    \section*{\listfigurename}%
    \@starttoc{lof}%
}
\newcommand{\customlistoftables}{%
    \section*{\listtablename}%
    \@starttoc{lot}%
}
\makeatother

% =====================================================
% HEADER / FOOTER
% =====================================================
\usepackage{fancyhdr}

\pagestyle{fancy}
\fancyhf{}
\fancyhead[L]{\small Báo cáo thực tập cơ sở}
\fancyhead[R]{\small \thepage}
\renewcommand{\headrulewidth}{0.4pt}

% =====================================================
% CÁC LỆNH TÙY CHỈNH
% =====================================================
\newcommand{\blankpage}{
    \newpage
    \thispagestyle{empty}
    \null
    \newpage
}

\newcommand{\coverpage}{
    \thispagestyle{empty}
    \pagenumbering{gobble}
}

\setcounter{tocdepth}{2}

% =====================================================
% BẮT ĐẦU TÀI LIỆU
% =====================================================
\begin{document}

% Đánh số hình, bảng, phương trình theo chương
\numberwithin{figure}{chapter}
\numberwithin{table}{chapter}
\numberwithin{equation}{chapter}


\setcounter{chapter}{3}

% =====================================================
% CHƯƠNG IV: XÂY DỰNG HỆ THỐNG
% =====================================================
\chapter{XÂY DỰNG HỆ THỐNG}

\section{Môi trường và công cụ phát triển}

\subsection{Nền tảng phần cứng}
Hệ thống Poliwise được triển khai trên kiến trúc microservices với Docker container. Môi trường phát triển và triển khai yêu cầu cấu hình phần cứng như sau:
\begin{itemize}
    \item \textbf{CPU:} Tối thiểu 4 cores (khuyến nghị 8 cores cho môi trường production).
    \item \textbf{RAM:} Tối thiểu 8 GB (khuyến nghị 16 GB).
    \item \textbf{Ổ cứng:} Tối thiểu 20 GB dung lượng trống cho image và các data volumes.
    \item \textbf{Hệ điều hành:} Linux (Ubuntu 22.04 LTS), macOS, hoặc Windows 10/11 với WSL2.
\end{itemize}

\subsection{Công nghệ và framework}

\subsubsection{Frontend}
\begin{table}[H]
\centering
\caption{Công nghệ sử dụng ở Frontend}
\label{tab:frontend_tech}
\begin{tabularx}{\textwidth}{|l|X|l|}
\hline
\rowcolor{gray!10}
Thành phần & Công nghệ & Phiên bản \\ \hline
Framework & Next.js (App Router) & 16.1.6 \\ \hline
Ngôn ngữ & TypeScript & 5.x \\ \hline
UI Library & Tailwind CSS & 4 \\ \hline
State Management & Zustand & 5.0.11 \\ \hline
HTTP Client & Axios & 1.13.6 \\ \hline
Data Visualization & Recharts & 3.8.1 \\ \hline
React Query & TanStack Query & 3.39.3 \\ \hline
Build Tool & Turbopack & Tích hợp Next.js \\ \hline
\end{tabularx}
\end{table}

\subsubsection{API Gateway}
\begin{table}[H]
\centering
\caption{Công nghệ sử dụng ở API Gateway}
\label{tab:gateway_tech}
\begin{tabularx}{\textwidth}{|l|X|l|}
\hline
\rowcolor{gray!10}
Thành phần & Công nghệ & Phiên bản \\ \hline
Framework & NestJS & 11.0.1 \\ \hline
Ngôn ngữ & TypeScript & 5.7.3 \\ \hline
Xác thực & Passport JWT & 4.0.1 \\ \hline
Rate Limiting & NestJS Throttler & 6.4.0 \\ \hline
Circuit Breaker & Opossum & 9.0.0 \\ \hline
Logging & Winston + Nest-Winston & 3.x / 1.10.x \\ \hline
Package Manager & pnpm & 9.15.4 \\ \hline
\end{tabularx}
\end{table}

\subsubsection{Microservices (Spring Boot)}
\begin{table}[H]
\centering
\caption{Công nghệ sử dụng ở các Microservices}
\label{tab:microservices_tech}
\begin{tabularx}{\textwidth}{|l|X|l|}
\hline
\rowcolor{gray!10}
Thành phần & Công nghệ & Phiên bản \\ \hline
Framework & Spring Boot & 3.3.4 (auth) / 3.4.3 (các service khác) \\ \hline
Ngôn ngữ & Java & 17 \\ \hline
Build Tool & Maven & 3.x (wrapper) \\ \hline
Database & PostgreSQL + pgvector & 16 \\ \hline
Message Queue & RabbitMQ & 3.13 \\ \hline
API Client & Spring Cloud OpenFeign & 2024.0.3 \\ \hline
Cache & Redis & 7 \\ \hline
Object Storage & MinIO (S3-compatible) & latest \\ \hline
Collaborative Editing & OnlyOffice Document Server & 8.1.0 \\ \hline
\end{tabularx}
\end{table}

\subsubsection{AI Services (Python)}
\begin{table}[H]
\centering
\caption{Công nghệ sử dụng ở các AI Services}
\label{tab:ai_services_tech}
\begin{tabularx}{\textwidth}{|l|X|l|}
\hline
\rowcolor{gray!10}
Thành phần & Công nghệ & Phiên bản \\ \hline
Framework & FastAPI & 0.115+ \\ \hline
Ngôn ngữ & Python & 3.11 \\ \hline
Embedding & HuggingFace TEI (BGE-M3) & CPU-1.6 \\ \hline
LLM & Groq API (Qwen3-8B / Gemini Flash) & - \\ \hline
\end{tabularx}
\end{table}

\subsection{Công cụ hỗ trợ phát triển}
\begin{table}[H]
\centering
\caption{Các công cụ hỗ trợ phát triển}
\label{tab:dev_tools}
\begin{tabularx}{\textwidth}{|l|X|}
\hline
\rowcolor{gray!10}
Mục đích & Công cụ \\ \hline
Container orchestration & Docker Compose \\ \hline
Database management & pgAdmin / DBeaver \\ \hline
API testing & Postman / curl \\ \hline
Version control & Git \\ \hline
Documentation & Markdown, Next.js MDX \\ \hline
Code quality & ESLint, Prettier \\ \hline
Monitoring & RabbitMQ Management UI (port 15672), MinIO Console (port 9001), Actuator endpoints \\ \hline
\end{tabularx}
\end{table}

\subsection{Cấu trúc thư mục project}
Cấu trúc thư mục chính của dự án được tổ chức như sau:
\begin{verbatim}
Poliwise/
├── frontend/web/             # Next.js 16 frontend
├── services/
│   ├── api-gateway/          # NestJS API gateway
│   ├── auth-service/          # Spring Boot - Authentication
│   ├── user-service/         # Spring Boot - User management
│   ├── knowledge-service/     # Spring Boot - Document management
│   ├── metadata-service/     # Spring Boot - Metadata/categories
│   ├── feedback-service/     # Spring Boot - Analytics/feedback
│   ├── ai-qa-service/       # FastAPI - AI Q&A
│   └── ingestion-service/   # FastAPI - Document processing
├── infrastructure/
│   └── init-db/             # SQL initialization scripts
├── config/                  # Shared environment config
├── contexts/                # Architecture documentation
├── docs/                    # Product notes
├── docker-compose.yml       # Full stack Docker orchestration
└── .env.example             # Environment template
\end{verbatim}

\section{Mô tả các chức năng chính}

\subsection{Xác thực và quản lý tài khoản (auth-service)}

\subsubsection{Đăng nhập / Đăng ký}
\begin{itemize}
    \item \textbf{Đăng nhập:} Người dùng nhập email và mật khẩu để xác thực. Hệ thống trả về JWT access token (15 phút) và refresh token (7 ngày). Tài khoản bị khóa tạm thời sau 5 lần đăng nhập thất bại trong 15 phút.
    \item \textbf{Đăng ký:} Admin tạo tài khoản mới với username, email, vai trò (\texttt{USER}/\texttt{MANAGER}/\texttt{ADMIN}), và phòng ban. Email được gửi kèm mật khẩu tạm.
    \item \textbf{Đăng ký hàng loạt:} Admin upload file CSV chứa danh sách người dùng để tạo nhiều tài khoản cùng lúc.
    \item \textbf{Đăng xuất:} Thu hồi refresh token hiện tại. Tùy chọn đăng xuất tất cả thiết bị.
\end{itemize}

\subsubsection{Quản lý phiên}
\begin{itemize}
    \item \textbf{Nhiều phiên đồng thời:} Mỗi lần đăng nhập tạo một phiên mới với thông tin thiết bị (device info), địa chỉ IP, và thời gian hết hạn.
    \item \textbf{Xem danh sách phiên:} Người dùng xem tất cả thiết bị đang đăng nhập.
    \item \textbf{Thu hồi phiên:} Đăng xuất từng phiên hoặc đăng xuất tất cả các phiên.
\end{itemize}

\subsubsection{Quản lý mật khẩu}
\begin{itemize}
    \item \textbf{Đổi mật khẩu:} Người dùng đổi mật khẩu (yêu cầu nhập mật khẩu hiện tại).
    \item \textbf{Yêu cầu mật khẩu:} Mật khẩu phải có độ dài tối thiểu 8 ký tự, bao gồm cả chữ hoa, chữ thường, chữ số và ký tự đặc biệt.
\end{itemize}

\subsubsection{Phân quyền RBAC}
Hệ thống phân quyền theo ba vai trò chính được thể hiện chi tiết tại Bảng \ref{tab:rbac_roles}:
\begin{table}[H]
\centering
\caption{Phân quyền theo vai trò trong hệ thống Poliwise}
\label{tab:rbac_roles}
\begin{tabularx}{\textwidth}{|l|X|}
\hline
\rowcolor{gray!10}
Vai trò & Khả năng hoạt động \\ \hline
\textbf{USER} & Hỏi đáp AI, quản lý hồ sơ cá nhân, gửi phản hồi. \\ \hline
\textbf{MANAGER} & Quyền của \texttt{USER} + xem dashboard phân tích, xuất báo cáo, xem danh sách câu hỏi chưa trả lời. \\ \hline
\textbf{ADMIN} & Quyền của \texttt{MANAGER} + tải tài liệu lên kho tri thức, quản lý metadata, quản lý người dùng, quản lý phòng ban. \\ \hline
\end{tabularx}
\end{table}

\subsection{Quản lý hồ sơ người dùng (user-service)}

\subsubsection{Thông tin cá nhân}
\begin{itemize}
    \item Người dùng có thể xem và chỉnh sửa các thông tin cá nhân bao gồm: họ tên, số điện thoại, chức vụ, giới thiệu bản thân, ngày sinh và phòng ban.
    \item Ảnh đại diện (avatar) của người dùng được tự động sinh ngẫu nhiên từ chữ cái đầu tiên của tên.
\end{itemize}

\subsubsection{Quản lý phòng ban}
\begin{itemize}
    \item \textbf{Admin tạo/cập nhật/xóa phòng ban:} Mỗi phòng ban bao gồm tên, mã phòng ban (viết hoa, duy nhất), mô tả và thông tin phòng ban cha (hỗ trợ cấu trúc cây phân cấp nhiều cấp).
    \item \textbf{Gán nhân viên:} Admin thực hiện gán người dùng vào phòng ban tương ứng, hệ thống sẽ hiển thị cảnh báo nếu chuyển người dùng từ phòng ban khác sang.
    \item \textbf{Xem nhân viên theo phòng ban:} Liệt kê danh sách người dùng thuộc từng phòng ban cụ thể.
\end{itemize}

\subsubsection{Trạng thái tài khoản}
Các trạng thái tài khoản trong hệ thống bao gồm:
\begin{itemize}
    \item \textbf{ACTIVE:} Tài khoản đang hoạt động bình thường.
    \item \textbf{DEACTIVATED:} Tài khoản bị vô hiệu hóa tạm thời (có thể kích hoạt lại bởi quản trị viên).
    \item \textbf{REVOKED:} Tài khoản bị thu hồi vĩnh viễn (không thể kích hoạt lại, hệ thống tự động khóa khi người dùng đổi mật khẩu thất bại quá số lần quy định).
    \item \textbf{Soft delete (Xóa mềm):} Xóa tài khoản nhưng không xóa vật lý khỏi cơ sở dữ liệu mà tiến hành ẩn danh (anonymize) các thông tin định danh cá nhân (PII).
\end{itemize}

\subsection{Hỏi đáp AI (ai-qa-service)}

\subsubsection{Chat AI}
\begin{itemize}
    \item Người dùng nhập câu hỏi liên quan đến chính sách nội bộ (quy trình, chính sách nhân sự, chế độ lương thưởng, tuyển dụng...).
    \item AI tổng hợp và trả lời dựa trên nội dung tài liệu đã được nạp (ingest) vào hệ thống.
    \item Phản hồi được truyền phát (streaming) theo thời gian thực (real-time token-by-token).
    \item Mỗi câu hỏi được liên kết với một cuộc trò chuyện (conversation) và lịch sử đàm thoại được lưu trữ đầy đủ.
\end{itemize}

\subsubsection{Nguồn tài liệu tham khảo}
\begin{itemize}
    \item Mỗi câu trả lời từ AI luôn đi kèm với danh sách các tài liệu nguồn đã được tham chiếu.
    \item Hiển thị điểm số mức độ liên quan (relevance score \%) đối với từng tài liệu.
    \item Trích xuất đoạn văn bản cụ thể được sử dụng làm cơ sở trả lời từ tài liệu gốc.
\end{itemize}

\subsubsection{Phản hồi và cải thiện}
\begin{itemize}
    \item Người dùng có thể đánh giá chất lượng câu trả lời bằng cách nhấn Like (thích) hoặc Dislike (không thích).
    \item Các câu hỏi mà AI trả lời không tốt hoặc người dùng đánh dấu chưa hài lòng sẽ được chuyển vào hàng đợi câu hỏi chưa trả lời để quản lý xem xét.
    \item Ghi nhận phản hồi tiêu cực để làm dữ liệu huấn luyện và tối ưu hóa hệ thống RAG sau này.
\end{itemize}

\subsubsection{Quản lý cuộc trò chuyện}
\begin{itemize}
    \item Lịch sử các cuộc trò chuyện được lưu trữ trên hệ thống.
    \item Người dùng có thể tìm kiếm hoặc xóa các cuộc trò chuyện cũ.
    \item Tiêu đề cuộc trò chuyện được AI tự động sinh dựa trên nội dung câu hỏi đầu tiên.
\end{itemize}

\subsection{Quản lý tài liệu (knowledge-service)}

\subsubsection{Upload tài liệu}
\begin{itemize}
    \item Quản trị viên tải lên tài liệu tri thức với các định dạng hỗ trợ: PDF, Word (DOCX/DOC), Excel (XLSX), Text, Markdown và File ảnh.
    \item Kích thước tệp tin tối đa được phép tải lên là 100MB.
    \item Quy trình tải lên gồm 4 bước trực quan: chọn tệp tin, theo dõi tiến trình tải lên, xác nhận metadata (tiêu đề, mô tả, danh mục phân loại, tags, ngôn ngữ, tùy chọn ``là chính sách'') và hoàn tất.
    \item Sau khi tải lên thành công, tài liệu sẽ ở trạng thái \texttt{STAGING} để chờ quản trị viên phê duyệt.
\end{itemize}

\subsubsection{Xử lý tài liệu (Ingestion Pipeline)}
Khi metadata của tài liệu được xác nhận, hệ thống sẽ tự động kích hoạt đường ống xử lý:
\begin{enumerate}
    \item \textbf{Parse (Trích xuất):} Trích xuất nội dung văn bản thuần túy từ tệp gốc (hỗ trợ OCR tự động cho tệp ảnh hoặc PDF quét).
    \item \textbf{Chunk (Phân mảnh):} Chia nội dung văn bản thành các phân mảnh nhỏ (chunks) phù hợp với giới hạn ngữ cảnh của mô hình nhúng.
    \item \textbf{Embed (Nhúng):} Chuyển đổi mỗi phân mảnh văn bản thành vector đặc trưng bằng mô hình BGE-M3 \cite{bgem32024}.
    \item \textbf{Index (Đánh chỉ mục):} Lưu trữ các vector và metadata tương ứng vào cơ sở dữ liệu vector (pgvector).
    \item \textbf{Trạng thái xử lý:} Di chuyển qua các trạng thái \texttt{STAGING} $\rightarrow$ \texttt{PARSING} $\rightarrow$ \texttt{READY} (hoặc \texttt{FAILED} nếu xảy ra lỗi trong quá trình xử lý).
\end{enumerate}

\subsubsection{Chỉnh sửa cộng tác (OnlyOffice)}
\begin{itemize}
    \item Tài liệu định dạng Word (DOCX/DOC) có thể được mở trực tiếp để soạn thảo bằng OnlyOffice ngay trên trình duyệt web.
    \item Hỗ trợ nhiều người dùng tham gia đồng biên tập theo thời gian thực.
    \item Cơ chế khóa tài liệu (document locking) giúp tránh xung đột: khi một người đang chỉnh sửa, những người dùng khác sẽ nhận được thông báo xung đột và chỉ có quyền xem.
    \item Sau khi lưu chỉnh sửa, phiên bản mới được tạo và hệ thống tự động thực hiện lại quá trình re-index cho tài liệu.
\end{itemize}

\subsubsection{Quản lý phiên bản}
\begin{itemize}
    \item Mỗi lượt tải lên tệp mới hoặc lưu thay đổi từ OnlyOffice sẽ tạo ra một phiên bản tài liệu mới.
    \item Người dùng có thể xem lịch sử phiên bản bao gồm: số phiên bản, ngày tạo, người cập nhật và ghi chú thay đổi.
    \item Hỗ trợ tải xuống phiên bản cũ, chỉnh sửa phiên bản cũ hoặc xóa bỏ (chỉ dành cho Admin).
    \item Tích hợp tính năng so sánh trực quan sự khác biệt giữa hai phiên bản tài liệu bất kỳ.
\end{itemize}

\subsubsection{Tìm kiếm và lọc}
\begin{itemize}
    \item Tìm kiếm toàn văn (full-text search) theo tên tệp tin hoặc nội dung bên trong tài liệu.
    \item Hỗ trợ bộ lọc nâng cao theo: loại tệp (PDF, Word, Excel, Text, Image), trạng thái tài liệu (\texttt{READY}, \texttt{STAGING}, \texttt{PARSING}, \texttt{FAILED}) và danh mục phân loại.
    \item Hỗ trợ 2 chế độ hiển thị: Dạng lưới (Grid view) và Dạng danh sách (List view).
    \item Hỗ trợ phân trang với cấu hình mặc định là 20 tài liệu trên mỗi trang.
\end{itemize}

\subsection{Metadata và phân loại (metadata-service)}

\subsubsection{Danh mục (Categories)}
\begin{itemize}
    \item Phân loại tài liệu theo cấu trúc hình cây (hỗ trợ danh mục cha - con).
    \item Admin tạo/sửa/xóa danh mục bao gồm các trường: tên danh mục, mô tả, biểu tượng (icon), danh mục cha và thứ tự hiển thị.
    \item Chuỗi slug thân thiện với đường dẫn URL được tự động sinh từ tên danh mục.
    \item Chỉ các danh mục ở trạng thái hoạt động (\texttt{active}) mới hiển thị trên giao diện người dùng.
\end{itemize}

\subsubsection{Tags (Nhãn dán)}
\begin{itemize}
    \item Cho phép gắn nhiều tags cho cùng một tài liệu để phân loại chéo.
    \item Admin tạo/sửa/xóa tag bao gồm: tên tag, màu sắc (16 màu cấu hình sẵn hoặc mã màu tùy chọn) và biểu tượng.
    \item Các tags phổ biến được làm nổi bật dựa trên tần suất sử dụng thực tế.
\end{itemize}

\subsubsection{Metadata tài liệu}
\begin{itemize}
    \item Mỗi tài liệu lưu trữ thông tin metadata đi kèm bao gồm: mô tả ngắn, ngôn ngữ sử dụng, trạng thái vòng đời tài liệu (\texttt{DRAFT}, \texttt{PUBLISHED}, \texttt{ARCHIVED}, \texttt{DELETED}).
    \item Metadata được khởi tạo ngay sau khi tài liệu được upload và xác nhận.
\end{itemize}

\subsubsection{Quy tắc truy cập (Access Rules)}
\begin{itemize}
    \item Quản trị viên có thể thiết lập chi tiết quyền truy cập tài liệu ở 3 cấp độ:
    \begin{itemize}
        \item \textbf{ROLE (Vai trò):} Cấp quyền theo vai trò tài khoản (\texttt{USER}, \texttt{MANAGER}, \texttt{ADMIN}).
        \item \textbf{DEPARTMENT (Phòng ban):} Cấp quyền theo đơn vị phòng ban của người dùng.
        \item \textbf{USER (Cá nhân):} Cấp quyền cụ thể cho từng tài khoản nhân viên.
    \end{itemize}
    \item Mỗi quy tắc quy định quyền truy cập dạng \texttt{VIEW} (Cho phép xem) hoặc \texttt{DENY} (Từ chối truy cập).
    \item \textbf{Mô phỏng truy cập:} Công cụ cho phép xem trước danh sách những người dùng được phép hoặc bị cấm xem tài liệu kèm theo lý do cụ thể để kiểm tra tính chính xác của các quy tắc.
\end{itemize}

\subsection{Phân tích và báo cáo (feedback-service)}

\subsubsection{Dashboard phân tích}
\begin{itemize}
    \item \textbf{6 chỉ số thống kê tổng quan:} Tổng số câu hỏi đã đặt, tỷ lệ hài lòng (like ratio), số tài liệu đang hoạt động, số người dùng tích cực, tình trạng hoạt động của API, tổng số lượt đăng nhập trong tháng.
    \item \textbf{Xu hướng sử dụng:} Biểu đồ hiển thị số lượng câu hỏi, số lượt like/dislike theo ngày/tuần/tháng.
    \item \textbf{Phân tích phản hồi:} Biểu đồ hình quạt phân bố tỷ lệ Like/Dislike của người dùng.
    \item \textbf{Top câu hỏi phổ biến:} Danh sách 5 câu hỏi được người dùng đặt nhiều nhất.
    \item \textbf{Top tài liệu trích dẫn:} Danh sách 5 tài liệu được hệ thống RAG trích dẫn nhiều nhất khi sinh câu trả lời.
    \item \textbf{Sức khỏe API:} Thống kê chi tiết số lượng request, tỷ lệ request thành công và thời gian phản hồi trung bình của từng endpoint.
\end{itemize}

\subsubsection{Báo cáo thống kê}
\begin{itemize}
    \item Hỗ trợ tạo và xuất 7 loại báo cáo nghiệp vụ: Tổng hợp sử dụng hệ thống, Phân tích câu hỏi người dùng, Phân tích phản hồi, Mức độ hoạt động của nhân viên, Mức độ phổ biến của tài liệu chính sách, Danh sách câu hỏi chưa được giải đáp, Thống kê hoạt động theo phòng ban.
    \item Định dạng tệp xuất ra: CSV hoặc JSON.
    \item Trạng thái tiến trình xuất báo cáo: \texttt{Pending} $\rightarrow$ \texttt{Processing} $\rightarrow$ \texttt{Completed} / \texttt{Failed}.
    \item Cho phép người dùng tải xuống các tệp báo cáo đã xuất thành công.
\end{itemize}

\subsubsection{Câu hỏi chưa được trả lời}
\begin{itemize}
    \item Thu thập các câu hỏi mà AI không tìm thấy tài liệu phù hợp hoặc bị người dùng đánh giá không hữu ích.
    \item Quản lý trạng thái câu hỏi: \texttt{PENDING} $\rightarrow$ \texttt{REVIEWING} $\rightarrow$ \texttt{ANSWERED} / \texttt{REJECTED}.
    \item Manager hoặc Admin có thể nhập câu trả lời thủ công cho câu hỏi này.
    \item Sau khi có câu trả lời thủ công, hệ thống sẽ lưu thông tin này để làm giàu kho tri thức phục vụ cho các câu hỏi tương tự sau đó.
\end{itemize}

\subsubsection{Nhật ký kiểm toán (Audit Logs)}
\begin{itemize}
    \item Ghi nhận tất cả các hành động của người dùng trên hệ thống: định danh người thực hiện, loại hành động, thời gian thực hiện và địa chỉ IP.
    \item Bộ lọc tìm kiếm lịch sử audit log dành cho Admin: lọc theo hành động, theo tài khoản người dùng và khoảng thời gian.
    \item Hỗ trợ phân trang để quản trị viên dễ dàng theo dõi và điều tra các sự cố bảo mật.
\end{itemize}

\subsection{API Gateway (api-gateway)}

Trong hệ thống PoliWise, API Gateway là cổng truy cập chính cho các API nghiệp vụ thông thường. Riêng thao tác upload tài liệu dạng \texttt{multipart/form-data} được xử lý như một ngoại lệ kỹ thuật: Frontend gửi trực tiếp đến \texttt{knowledge-service} để đảm bảo truyền file lớn ổn định, trong khi JWT và quyền truy cập vẫn được kiểm tra tại service đích.

\subsubsection{JWT Authentication}
\begin{itemize}
    \item Kiểm tra tính hợp lệ của JWT access token truyền qua HTTP Header \texttt{Authorization: Bearer <token>}.
    \item Xác minh trạng thái tài khoản người dùng (\texttt{ACTIVE}/\texttt{DEACTIVATED}/\texttt{REVOKED}).
    \item Hỗ trợ cơ chế tự động gia hạn token (token refresh) khi access token sắp hết hạn.
\end{itemize}

\subsubsection{Rate Limiting (Giới hạn tần suất)}
Giới hạn tần suất được thiết lập chi tiết theo từng nhóm vai trò như trong Bảng \ref{tab:rate_limiting}:
\begin{table}[H]
\centering
\caption{Cấu hình giới hạn tần suất yêu cầu tại API Gateway}
\label{tab:rate_limiting}
\begin{tabularx}{\textwidth}{|l|X|}
\hline
\rowcolor{gray!10}
Vai trò & Giới hạn yêu cầu tối đa \\ \hline
\texttt{PUBLIC} (Chưa đăng nhập) & 200 requests / phút \\ \hline
\texttt{USER} & 100 requests / phút \\ \hline
\texttt{MANAGER} & 200 requests / phút \\ \hline
\texttt{ADMIN} & 500 requests / phút \\ \hline
\end{tabularx}
\end{table}

\subsubsection{Circuit Breaker (Bộ ngắt mạch)}
\begin{itemize}
    \item Giúp ngăn chặn sự cố tràn lan (cascade failure) giữa các dịch vụ trong hệ thống microservices.
    \item Khi một service downstream gặp sự cố hoặc quá tải, Circuit Breaker sẽ tự động mở mạch và trả về phản hồi mặc định (fallback response) thay vì tiếp tục gửi yêu cầu làm nghẽn hệ thống.
\end{itemize}

\subsubsection{Health Checks}
\begin{itemize}
    \item \texttt{/health}: Kiểm tra tình trạng sức khỏe tổng thể của toàn bộ hệ thống.
    \item \texttt{/health/live}: Liveness probe dùng để xác định API Gateway đang chạy.
    \item \texttt{/health/ready}: Readiness probe dùng để xác nhận toàn bộ các service phụ trợ đã sẵn sàng kết nối.
    \item \texttt{/health/circuit-breakers}: Kiểm tra trạng thái đóng/mở của các bộ ngắt mạch.
    \item \texttt{/health/api-metrics}: Thống kê các số liệu hiệu năng của API (số lượng request, latency, error rate).
\end{itemize}

\subsection{Giao diện người dùng (frontend/web)}
Giao diện Next.js sử dụng cấu trúc thư mục App Router mới. Các màn hình giao diện chính và phân quyền truy cập tương ứng được liệt kê trong Bảng \ref{tab:routes_list}:

\begin{table}[H]
\centering
\caption{Danh sách các màn hình giao diện trên Frontend}
\label{tab:routes_list}
\begin{tabularx}{\textwidth}{|l|X|l|}
\hline
\rowcolor{gray!10}
Đường dẫn (Route) & Nội dung mô tả trang & Vai trò tối thiểu \\ \hline
\texttt{/login} & Màn hình đăng nhập hệ thống & Public \\ \hline
\texttt{/register} & Màn hình đăng ký tài khoản & Public \\ \hline
\texttt{/reset-password} & Màn hình đặt lại mật khẩu mới & Public \\ \hline
\texttt{/} & Giao diện Chat AI chính & \texttt{USER} \\ \hline
\texttt{/documents} & Thư viện quản lý tài liệu tri thức & \texttt{USER} \\ \hline
\texttt{/documents/[id]} & Chi tiết tài liệu (6 tab thông tin và cấu hình) & \texttt{USER} \\ \hline
\texttt{/documents/compare} & Công cụ so sánh các phiên bản tài liệu & \texttt{USER} \\ \hline
\texttt{/profile} & Hồ sơ cá nhân (thông tin, chỉnh sửa) & \texttt{USER} \\ \hline
\texttt{/sessions} & Quản lý danh sách các thiết bị đang đăng nhập & \texttt{USER} \\ \hline
\texttt{/analytics} & Dashboard hiển thị các chỉ số phân tích dữ liệu & \texttt{MANAGER} \\ \hline
\texttt{/analytics/reports} & Màn hình yêu cầu và tải xuống báo cáo thống kê & \texttt{MANAGER} \\ \hline
\texttt{/admin/users} & Quản trị tài khoản người dùng & \texttt{ADMIN} \\ \hline
\texttt{/admin/departments} & Quản trị cơ cấu tổ chức và sơ đồ phòng ban & \texttt{ADMIN} \\ \hline
\texttt{/admin/metadata/categories} & Quản lý danh mục tài liệu phân cấp & \texttt{ADMIN} \\ \hline
\texttt{/admin/metadata/tags} & Quản lý danh sách các thẻ tags & \texttt{ADMIN} \\ \hline
\texttt{/admin/unanswered} & Xem và xử lý các câu hỏi chưa được giải đáp & \texttt{MANAGER} \\ \hline
\end{tabularx}
\end{table}

\subsection{Tóm tắt luồng nghiệp vụ chính}

\subsubsection{Luồng 1: Hỏi đáp AI}
\begin{verbatim}
Người dùng nhập câu hỏi
  → Frontend gửi sang API Gateway
    → API Gateway validate JWT, check rate limit
      → Proxy sang ai-qa-service
        → Layer 1: lọc an toàn đầu vào
          → Layer 2: phân loại SIMPLE / COMPLEX
            → Nếu SIMPLE: trả lời trực tiếp bằng model nhanh
            → Nếu COMPLEX:
              → Query Refiner tối ưu câu hỏi cho tìm kiếm
              → Hybrid Search (BGE-M3 + BM25 + RRF)
              → Knowledge Gap Detection
              → Model Registry chọn model sinh câu trả lời
              → Trả lời streaming + sources về frontend
              → Lưu conversation/message vào database
              → Publish unanswered.question nếu phát hiện knowledge gap
\end{verbatim}

\begin{figure}[H]
\centering
\includegraphics[height=0.82\textheight,keepaspectratio]{images/ai-qa-pipeline-current.png}
\caption{Luồng xử lý hiện tại của AI Q\&A Service}
\label{fig:ai_qa_pipeline_current}
\end{figure}

\subsubsection{Lý do lựa chọn các mô hình trong AI Q\&A Service}

Hệ thống sử dụng nhiều mô hình khác nhau theo nguyên tắc phân tầng, không dùng một mô hình lớn cho toàn bộ pipeline. Các tác vụ nhẹ như lọc an toàn, phân loại ý định, viết lại câu hỏi và trả lời hội thoại đơn giản được xử lý bằng các mô hình nhanh để giảm độ trễ và tiết kiệm quota API. Những câu hỏi cần truy xuất tài liệu mới được chuyển sang lớp RAG để dùng embedding, tìm kiếm lai và mô hình sinh câu trả lời cuối cùng.

Ở Layer 1, hệ thống ưu tiên bộ lọc keyword/regex để chặn nhanh các mẫu toxic hoặc prompt injection rõ ràng. Với câu hỏi tiếng Anh, hệ thống dùng \texttt{meta-llama/llama-prompt-guard-2-86m} vì đây là mô hình nhỏ, chuyên biệt cho bài toán kiểm tra an toàn prompt. Với câu hỏi tiếng Việt, hệ thống dùng \texttt{llama-3.1-8b-instant} để đánh giá an toàn theo ngữ cảnh do mô hình này xử lý tiếng Việt tốt hơn, giúp giảm nguy cơ chặn nhầm các câu hỏi hợp lệ.

Ở Layer 2, \texttt{llama-3.1-8b-instant} được dùng cho phân loại intent, trả lời câu hỏi đơn giản, viết lại truy vấn và sinh tiêu đề hội thoại. Các tác vụ này thường có đầu ra ngắn, yêu cầu phản hồi nhanh và không cần truy xuất tài liệu. Vì vậy, dùng một mô hình nhanh cho các bước này phù hợp hơn so với gọi mô hình sinh câu trả lời lớn ngay từ đầu.

Ở Layer 3, hệ thống dùng \texttt{BAAI/bge-m3} thông qua HuggingFace Text Embeddings Inference để tạo embedding cho tìm kiếm ngữ nghĩa. BGE-M3 phù hợp với bài toán truy xuất tài liệu vì hỗ trợ đa ngôn ngữ và tạo vector 1024 chiều. Kết quả vector search được kết hợp với BM25 và Reciprocal Rank Fusion để xử lý tốt cả câu hỏi ngữ nghĩa lẫn câu hỏi chứa từ khóa, mã điều khoản hoặc thuật ngữ chính xác.

Đối với bước sinh câu trả lời cuối cùng, hệ thống sử dụng Model Registry. Mặc định là \texttt{local/qwen3-8b} nhằm ưu tiên khả năng tự host và kiểm soát dữ liệu. Khi cần chất lượng cao hơn, khi người dùng chọn model khác, hoặc khi model local không khả dụng, hệ thống có thể chuyển sang các model remote đã cấu hình như Groq Qwen3-32B, Groq Llama 3.3 70B, Gemini Flash hoặc OpenRouter Mistral tùy API key và trạng thái khả dụng của từng model. Ngoài luồng hỏi đáp, \texttt{ingestion-service} còn dùng \texttt{llama-3.3-70b-versatile} qua Groq cho tác vụ gợi ý metadata tài liệu vì tác vụ này cần đầu ra JSON có cấu trúc và không nằm trên đường phản hồi trực tiếp của chat.

\subsubsection{Luồng 2: Upload tài liệu}
\begin{verbatim}
Admin upload file
  → Frontend → knowledge-service (multipart upload, bypass Gateway kỹ thuật)
    → Lưu file vào MinIO
      → Publish event "ingestion.requested" qua RabbitMQ
        → ingestion-service nhận event
          → Parse file → Chunk → Embed → Index (pgvector)
            → Publish event "document.uploaded" và cập nhật trạng thái thành READY
              → Trả về cho admin xác nhận metadata
                → Admin xác nhận
                  → Document sẵn sàng cho AI Q&A
\end{verbatim}

\subsubsection{Luồng 3: Quản lý truy cập}
\begin{verbatim}
Admin tạo Access Rule cho tài liệu
  → Frontend → metadata-service
    → Lưu rule (ROLE/DEPARTMENT/USER + VIEW/DENY)
    → Tự động tạo metadata record nếu chưa có
      → Khi người dùng truy cập tài liệu
        → knowledge-service gọi metadata-service check permission
          → Trả về kết quả: cho phép hoặc từ chối
\end{verbatim}

\section{Giao diện minh họa}
Danh sách các tệp hình ảnh tương ứng với các màn hình giao diện của hệ thống Poliwise được quản lý tại thư mục \texttt{frontend/docs/images/} như sau:

\begin{enumerate}
    \item \texttt{/login} -- Màn hình đăng nhập
    \begin{figure}[H]
    \centering
    \includegraphics[width=0.85\textwidth]{images/04-login-page.png}
    \caption{Màn hình đăng nhập}
    \label{fig:04_login_page}
    \end{figure}
    
    \item \texttt{/register} -- Màn hình đăng ký
    \begin{figure}[H]
    \centering
    \includegraphics[width=0.85\textwidth]{images/04-register-page.png}
    \caption{Màn hình đăng ký (Giao diện nhập thông tin)}
    \label{fig:04_register_page}
    \end{figure}
    \begin{figure}[H]
    \centering
    \includegraphics[width=0.85\textwidth]{images/04-register-page-02.png}
    \caption{Màn hình đăng ký (Giao diện hoàn tất)}
    \label{fig:04_register_page_02}
    \end{figure}

    \item \texttt{/reset-password} -- Màn hình đặt lại mật khẩu
    \begin{figure}[H]
    \centering
    \includegraphics[width=0.85\textwidth]{images/04-reset-password-page.png}
    \caption{Màn hình đặt lại mật khẩu}
    \label{fig:04_reset_password_page}
    \end{figure}

    \item \texttt{/} -- Màn hình Chat AI (màn hình chào)
    \begin{figure}[H]
    \centering
    \includegraphics[width=0.85\textwidth]{images/04-home-chat-welcome.png}
    \caption{Màn hình Chat AI (màn hình chào)}
    \label{fig:04_home_chat_welcome}
    \end{figure}

    \item \texttt{/} -- Màn hình Chat AI (cuộc trò chuyện)
    \begin{figure}[H]
    \centering
    \includegraphics[width=0.85\textwidth]{images/04-home-chat-conversation.png}
    \caption{Màn hình Chat AI (cuộc trò chuyện)}
    \label{fig:04_home_chat_conversation}
    \end{figure}

    \item \texttt{/documents} -- Thư viện tài liệu (Dạng lưới)
    \begin{figure}[H]
    \centering
    \includegraphics[width=0.85\textwidth]{images/04-documents-grid.png}
    \caption{Thư viện tài liệu (Hiển thị dạng lưới)}
    \label{fig:04_documents_grid}
    \end{figure}

    \item \texttt{/documents} -- Thư viện tài liệu (Dạng danh sách)
    \begin{figure}[H]
    \centering
    \includegraphics[width=0.85\textwidth]{images/04-documents-list.png}
    \caption{Thư viện tài liệu (Hiển thị dạng danh sách)}
    \label{fig:04_documents_list}
    \end{figure}

    \item \texttt{/documents} -- Tiến trình 4 bước tải lên tài liệu
    \begin{figure}[H]
    \centering
    \includegraphics[width=0.85\textwidth]{images/04-documents-upload-wizard-01.png}
    \caption{Tiến trình tải lên tài liệu -- Bước 1: Chọn tệp tin}
    \label{fig:04_upload_wizard_01}
    \end{figure}
    \begin{figure}[H]
    \centering
    \includegraphics[width=0.85\textwidth]{images/04-documents-upload-wizard-02.png}
    \caption{Tiến trình tải lên tài liệu -- Bước 2: Nhập thông tin metadata}
    \label{fig:04_upload_wizard_02}
    \end{figure}

    \item \texttt{/documents/[id]} -- Các tab chi tiết tài liệu
    \begin{figure}[H]
    \centering
    \includegraphics[width=0.85\textwidth]{images/04-document-detail-tabs.png}
    \caption{Chi tiết tài liệu với các tab thông tin và cấu hình}
    \label{fig:04_document_detail_tabs}
    \end{figure}

    \item \texttt{/documents/[id]} -- Màn hình mô phỏng quyền truy cập
    \begin{figure}[H]
    \centering
    \includegraphics[width=0.85\textwidth]{images/04-document-access-simulation.png}
    \caption{Giao diện mô phỏng và kiểm tra quyền truy cập của người dùng}
    \label{fig:04_access_simulation}
    \end{figure}

    \item \texttt{/documents/compare} -- Công cụ so sánh tài liệu
    \begin{figure}[H]
    \centering
    \includegraphics[width=0.85\textwidth]{images/04-document-compare.png}
    \caption{Công cụ so sánh sự khác biệt giữa các phiên bản tài liệu}
    \label{fig:04_document_compare}
    \end{figure}

    \item \texttt{/profile} -- Hồ sơ cá nhân - Tab thông tin
    \begin{figure}[H]
    \centering
    \includegraphics[width=0.85\textwidth]{images/04-profile-info.png}
    \caption{Hồ sơ cá nhân -- Tab xem thông tin chi tiết}
    \label{fig:04_profile_info}
    \end{figure}

    \item \texttt{/profile} -- Hồ sơ cá nhân - Tab chỉnh sửa
    \begin{figure}[H]
    \centering
    \includegraphics[width=0.85\textwidth]{images/04-profile-edit.png}
    \caption{Hồ sơ cá nhân -- Tab thay đổi và cập nhật thông tin}
    \label{fig:04_profile_edit}
    \end{figure}

    \item \texttt{/sessions} -- Màn hình quản lý các phiên đăng nhập
    \begin{figure}[H]
    \centering
    \includegraphics[width=0.85\textwidth]{images/04-sessions-page.png}
    \caption{Màn hình quản lý danh sách thiết bị và các phiên đăng nhập}
    \label{fig:04_sessions_page}
    \end{figure}

    \item \texttt{/analytics} -- Dashboard thống kê phân tích hoạt động
    \begin{figure}[H]
    \centering
    \includegraphics[width=0.85\textwidth]{images/04-analytics-dashboard.png}
    \caption{Dashboard thống kê và phân tích tần suất sử dụng hệ thống}
    \label{fig:04_analytics_dashboard}
    \end{figure}

    \item \texttt{/analytics/reports} -- Màn hình tạo và tải báo cáo
    \begin{figure}[H]
    \centering
    \includegraphics[width=0.85\textwidth]{images/04-reports-page.png}
    \caption{Giao diện yêu cầu kết xuất và tải về báo cáo thống kê}
    \label{fig:04_reports_page}
    \end{figure}

    \item \texttt{/admin/users} -- Quản lý người dùng của hệ thống
    \begin{figure}[H]
    \centering
    \includegraphics[width=0.85\textwidth]{images/04-admin-users.png}
    \caption{Giao diện quản trị tài khoản người dùng và phân quyền hệ thống}
    \label{fig:04_admin_users}
    \end{figure}

    \item \texttt{/admin/departments} -- Quản lý cơ cấu phòng ban
    \begin{figure}[H]
    \centering
    \includegraphics[width=0.85\textwidth]{images/04-admin-departments.png}
    \caption{Giao diện quản lý sơ đồ và cơ cấu phòng ban tổ chức}
    \label{fig:04_admin_departments}
    \end{figure}

    \item \texttt{/admin/metadata/categories} -- Quản trị danh mục tài liệu
    \begin{figure}[H]
    \centering
    \includegraphics[width=0.85\textwidth]{images/04-admin-categories.png}
    \caption{Giao diện thiết lập và phân cấp danh mục tài liệu tri thức}
    \label{fig:04_admin_categories}
    \end{figure}

    \item \texttt{/admin/metadata/tags} -- Quản trị nhãn dán tài liệu
    \begin{figure}[H]
    \centering
    \includegraphics[width=0.85\textwidth]{images/04-admin-tags.png}
    \caption{Giao diện quản lý các nhãn dán (tags) tài liệu}
    \label{fig:04_admin_tags}
    \end{figure}

    \item \texttt{/admin/unanswered} -- Xử lý câu hỏi chưa có câu trả lời
    \begin{figure}[H]
    \centering
    \includegraphics[width=0.85\textwidth]{images/04-admin-unanswered.png}
    \caption{Giao diện xử lý câu hỏi chưa có câu trả lời}
    \label{fig:04_admin_unanswered}
    \end{figure}
\end{enumerate}

\section{Cài đặt và cấu hình hệ thống}

\subsection{Yêu cầu hệ thống}
Yêu cầu tối thiểu và khuyến nghị để triển khai hệ thống được tóm tắt trong Bảng \ref{tab:req_specs}:
\begin{table}[H]
\centering
\caption{Yêu cầu về hạ tầng phần cứng và phần mềm}
\label{tab:req_specs}
\begin{tabularx}{\textwidth}{|l|X|X|}
\hline
\rowcolor{gray!10}
Thành phần & Yêu cầu tối thiểu & Yêu cầu khuyến nghị \\ \hline
Bộ vi xử lý (CPU) & 4 cores & 8 cores \\ \hline
Bộ nhớ (RAM) & 8 GB & 16 GB \\ \hline
Không gian đĩa (Disk) & 20 GB & 50 GB \\ \hline
Phiên bản Docker & 20.10+ & 24.x \\ \hline
Phiên bản Docker Compose & 2.0+ & 2.20+ \\ \hline
\end{tabularx}
\end{table}

\subsection{Cài đặt môi trường phát triển}

\subsubsection{Bước 1: Tải mã nguồn dự án}
Mã nguồn dự án được lưu trữ tại GitHub: \url{https://github.com/Poliwise/Poliwise}. Sử dụng công cụ Git để clone mã nguồn của dự án về máy tính:
\begin{verbatim}
git clone https://github.com/Poliwise/Poliwise.git
cd Poliwise
\end{verbatim}

\subsubsection{Bước 2: Cấu hình tệp tin môi trường}
Sao chép tệp tin cấu hình mẫu \texttt{.env.example} thành \texttt{.env} và cập nhật các tham số bảo mật và thông tin kết nối:
\begin{verbatim}
cp .env.example .env
\end{verbatim}
Các tham số cấu hình chính cần được thiết lập:
\begin{verbatim}
# JWT - Tạo chuỗi ngẫu nhiên 64 ký tự hex làm khóa bí mật
JWT_SECRET=<generate-random-64-char-hex>

# Groq API Key - Sử dụng dịch vụ LLM đám mây
GROQ_API_KEY=<your-groq-api-key>

# Cơ sở dữ liệu PostgreSQL
POSTGRES_PASSWORD=<secure-password>

# Hàng đợi thông điệp RabbitMQ
RABBITMQ_PASSWORD=<secure-password>

# Lưu trữ đối tượng MinIO
MINIO_ROOT_PASSWORD=<secure-password>
\end{verbatim}

\subsubsection{Bước 3: Khởi động hệ thống thông qua Docker}
Chạy lệnh sau ở thư mục gốc để khởi chạy tất cả các dịch vụ:
\begin{verbatim}
docker compose up -d
\end{verbatim}
Kiểm tra danh sách các container đang chạy và đảm bảo trạng thái hoạt động tốt:
\begin{verbatim}
docker compose ps
\end{verbatim}

\subsubsection{Bước 4: Các địa chỉ truy cập dịch vụ}
Sau khi khởi chạy thành công, có thể truy cập các thành phần của hệ thống thông qua các địa chỉ trong Bảng \ref{tab:port_mappings}:
\begin{table}[H]
\centering
\caption{Bảng ánh xạ địa chỉ và cổng truy cập của các thành phần}
\label{tab:port_mappings}
\begin{tabularx}{\textwidth}{|l|X|}
\hline
\rowcolor{gray!10}
Thành phần dịch vụ & Địa chỉ truy cập \\ \hline
Giao diện người dùng (Frontend) & \url{http://localhost:3000} \\ \hline
Cổng API Gateway & \url{http://localhost:3001} \\ \hline
Trang quản lý RabbitMQ & \url{http://localhost:15672} \\ \hline
Trang quản lý MinIO Console & \url{http://localhost:9001} \\ \hline
OnlyOffice Document Server & \url{http://localhost:8888} \\ \hline
\end{tabularx}
\end{table}

\subsubsection{Bước 5: Tài khoản quản trị mặc định}
Hệ thống tạo sẵn một tài khoản quản trị cao cấp phục vụ cho lần đăng nhập đầu tiên:
\begin{verbatim}
Email đăng nhập: admin@poliwise.com
Mật khẩu mặc định: Admin@123456
Vai trò: ADMIN
\end{verbatim}

\subsection{Cài đặt các thành phần độc lập (không sử dụng Docker)}
Trong trường hợp cần phát triển hoặc debug riêng lẻ từng thành phần trực tiếp trên máy local:

\subsubsection{Frontend}
\begin{verbatim}
cd frontend/web
pnpm install
pnpm dev
# Giao diện chạy tại địa chỉ http://localhost:3000
\end{verbatim}

\subsubsection{API Gateway}
\begin{verbatim}
cd services/api-gateway
pnpm install
pnpm run start:dev
# API Gateway lắng nghe tại http://localhost:3001
\end{verbatim}

\subsubsection{Các Microservices nghiệp vụ (Spring Boot)}
\begin{verbatim}
cd services/<service-name>
./mvnw spring-boot:run
\end{verbatim}
Các cổng mặc định tương ứng của các dịch vụ Java bao gồm: 8081 (\texttt{auth-service}), 8082 (\texttt{user-service}), 8083 (\texttt{knowledge-service}), 8084 (\texttt{metadata-service}) và 8085 (\texttt{feedback-service}).

\subsection{Cấu hình RabbitMQ}
RabbitMQ hoạt động như một trục thông điệp trung tâm điều phối các sự kiện bất đồng bộ giữa các microservices, sử dụng cấu hình trao đổi như ở Bảng \ref{tab:rabbitmq_exchanges}:
\begin{table}[H]
\centering
\caption{Cấu hình các hàng đợi và sự kiện trong RabbitMQ}
\label{tab:rabbitmq_exchanges}
\begin{tabularx}{\textwidth}{|l|l|l|l|}
\hline
\rowcolor{gray!10}
Exchange Name & Exchange Type & Queue Name & Routing Key \\ \hline
\texttt{user.events} & topic & \texttt{user.status.changes} & \texttt{user.status.changed} \\ \hline
\texttt{user.events} & topic & \texttt{user.revoked} & \texttt{user.revoked} \\ \hline
\texttt{document.events} & topic & \texttt{document.uploaded} & \texttt{document.uploaded} \\ \hline
\texttt{document.events} & topic & \texttt{document.deleted} & \texttt{document.deleted} \\ \hline
\texttt{feedback.events} & topic & \texttt{unanswered.question} & \texttt{unanswered.question} \\ \hline
\end{tabularx}
\end{table}

\subsection{Cấu hình AI Services}

\subsubsection{AI Q\&A Service (ai-qa-service)}
Thiết lập các biến môi trường trong tệp \texttt{.env} của AI service để kết nối dịch vụ LLM:
\begin{verbatim}
LLM_BASE_URL=https://api.groq.com/openai/v1
GROQ_API_KEY=<your-api-key>
LOCAL_MODEL_NAME=qwen3-8b
GATEWAY_API_KEY=<gemini-api-key>
GATEWAY_MODEL=gemini-flash-latest
\end{verbatim}

\subsubsection{Embedding Service (BGE-M3)}
Chạy dịch vụ trích xuất vector embedding cục bộ thông qua Hugging Face Text Embeddings Inference (TEI):
\begin{verbatim}
docker run -p 8001:80 \
  -e MODEL_ID=BAAI/bge-m3 \
  ghcr.io/huggingface/text-embeddings-inference:cpu-1.6
\end{verbatim}

\subsection{Cấu hình OnlyOffice}
Kích hoạt mã bảo mật JWT secret trên OnlyOffice Document Server để kiểm soát quyền biên tập tài liệu:
\begin{verbatim}
docker exec poliwise-onlyoffice documentserver jumper \
  jwt_secret=<your-jwt-secret>
\end{verbatim}
Đồng thời, đảm bảo cấu hình chính xác biến môi trường \texttt{NEXT\_PUBLIC\_ONLYOFFICE\_URL} ở Frontend trỏ về địa chỉ IP của OnlyOffice container.

\subsection{Cấu hình MinIO}
MinIO đóng vai trò lưu trữ các tệp gốc tải lên của hệ thống:
\begin{itemize}
    \item Trang quản trị Web Console: \url{http://localhost:9001}
    \item S3 Endpoint URL: \url{http://localhost:9000}
    \item Bucket lưu trữ chính mặc định: \texttt{poliwise-documents}
\end{itemize}

\subsection{Kiểm tra trạng thái hệ thống}
Sử dụng các công cụ dòng lệnh hoặc trình duyệt để kiểm tra trạng thái hoạt động và hiệu năng của hệ thống:
\begin{verbatim}
# Kiểm tra sức khỏe tổng thể hệ thống
curl http://localhost:3001/health

# Kiểm tra tính sẵn sàng của các dịch vụ downstream
curl http://localhost:3001/health/ready
curl http://localhost:3001/health/circuit-breakers

# Lấy các chỉ số thống kê hiệu năng API
curl -H "Authorization: Bearer <token>" \
  http://localhost:3001/health/api-metrics
\end{verbatim}

\subsection{Database schema initialization}
Các kịch bản SQL khởi tạo cơ sở dữ liệu được lưu trữ tại thư mục \texttt{infrastructure/init-db/}. Các schema được phân bổ và tự động import khi PostgreSQL khởi chạy lần đầu:
\begin{table}[H]
\centering
\caption{Phân chia các database schema theo dịch vụ}
\label{tab:db_init_schemas}
\begin{tabularx}{\textwidth}{|l|X|}
\hline
\rowcolor{gray!10}
Tên Database Schema & Dịch vụ sở hữu trực tiếp \\ \hline
\texttt{poliwise\_auth} & auth-service \\ \hline
\texttt{poliwise\_user} & user-service \\ \hline
\texttt{poliwise\_knowledge} & knowledge-service \\ \hline
\texttt{poliwise\_metadata} & metadata-service \\ \hline
\texttt{poliwise\_analytics} & feedback-service \\ \hline
\end{tabularx}
\end{table}

\subsection{Dừng và dọn dẹp hệ thống}
Các lệnh Docker để quản lý và giải phóng tài nguyên hệ thống:
\begin{verbatim}
# Tạm dừng hoạt động của các container (dữ liệu trong volume được bảo toàn)
docker compose stop

# Dừng và xóa bỏ hoàn toàn container (dữ liệu trong volume được bảo toàn)
docker compose down

# Dừng và xóa bỏ container cùng toàn bộ dữ liệu trong volume
docker compose down -v
\end{verbatim}

\clearpage
% =====================================================
% CHƯƠNG V: ĐÁNH GIÁ KẾT QUẢ VÀ RÚT KINH NGHIỆM
% =====================================================
\chapter{ĐÁNH GIÁ KẾT QUẢ VÀ RÚT KINH NGHIỆM}

Chương này tổng hợp và phân tích toàn diện kết quả đạt được của hệ thống PoliWise sau quá trình triển khai thực nghiệm, so sánh hiệu quả với phương thức tra cứu truyền thống, đánh giá mức độ đáp ứng các mục tiêu ban đầu và đúc kết các bài học kinh nghiệm quý giá trong suốt quá trình nghiên cứu và xây dựng hệ thống.

\rule{\textwidth}{0.5pt}

\section{Đánh giá mức độ đạt mục tiêu đề tài}

\subsection{Đối chiếu kết quả với mục tiêu ban đầu}

Nhóm nghiên cứu tiến hành đối chiếu từng mục tiêu được đề ra tại Chương I với kết quả thực tế đạt được sau quá trình xây dựng và thử nghiệm hệ thống:

\begin{table}[H]
\centering
\small
\caption{Đối chiếu mục tiêu và kết quả đạt được}
\label{tab:goal_vs_result}
\begin{tabularx}{\textwidth}{|l|X|X|l|}
\hline
\rowcolor{gray!10}
STT & Mục tiêu đề ra & Kết quả đạt được & Trạng thái \\ \hline
1 & Hỏi đáp bằng ngôn ngữ tự nhiên & Chatbot phản hồi streaming, hỗ trợ tiếng Việt đầy đủ & Đạt \\ \hline
2 & Tìm kiếm ngữ nghĩa (Semantic Search) & Hybrid Search (Vector + BM25) + RRF, có cơ chế Reranking cấu hình được, đạt Context Precision 89.5\% trong bộ thử nghiệm nội bộ & Đạt \\ \hline
3 & Trích dẫn nguồn minh bạch & Mỗi câu trả lời kèm đoạn trích dẫn và metadata tài liệu gốc & Đạt \\ \hline
4 & Phân quyền RBAC theo vai trò & Kiểm soát truy cập theo Role + Department ngay tại tầng vector search & Đạt \\ \hline
5 & So sánh phiên bản tài liệu & Giao diện diff 3 chiều hỗ trợ phát hiện thay đổi giữa các phiên bản & Đạt \\ \hline
6 & Ghi nhận phản hồi người dùng & Module Like/Dislike + Comment + Analytics Dashboard đã hoạt động & Đạt \\ \hline
7 & Triển khai On-premise & Toàn bộ hệ thống chạy trong Docker Compose, không phụ thuộc Cloud & Đạt \\ \hline
\end{tabularx}
\end{table}

Kết quả cho thấy hệ thống PoliWise đã hoàn thành các mục tiêu chức năng chính đề ra ban đầu. Một số tiêu chí định lượng như độ trung thực câu trả lời (Faithfulness 93.8\%) và thời gian phản hồi (End-to-End Latency 1.2 giây) được ghi nhận trong phạm vi thử nghiệm nội bộ quy mô nhỏ, dùng để tham khảo và cần tiếp tục mở rộng kiểm chứng khi triển khai thực tế.

\subsection{Đánh giá chất lượng RAG theo khung Ragas}

Dựa trên kết quả thực nghiệm nội bộ với bộ 150 câu hỏi kiểm thử mô tả ở Chương IV, bảng sau tóm tắt các chỉ số chất lượng RAG thông qua khung đánh giá Ragas \cite{ragas2024}. Các số liệu này phản ánh tập dữ liệu thử nghiệm của nhóm, chưa đại diện cho toàn bộ môi trường doanh nghiệp thực tế:

\begin{table}[H]
\centering
\caption{Tổng hợp chỉ số chất lượng RAG (Ragas Framework)}
\label{tab:ragas_summary}
\begin{tabularx}{\textwidth}{|l|X|c|l|}
\hline
\rowcolor{gray!10}
Chỉ số & Ý nghĩa & Điểm đạt & Nhận xét \\ \hline
Context Precision & Tỷ lệ chunk truy xuất thực sự liên quan & 89.5\% & Tốt \\ \hline
Context Recall & Tỷ lệ thông tin quan trọng không bị bỏ sót & 91.2\% & Tốt \\ \hline
Faithfulness & Câu trả lời dựa trên tài liệu, không bịa đặt & 93.8\% & Xuất sắc \\ \hline
Answer Relevance & Câu trả lời bám sát câu hỏi của người dùng & 88.2\% & Tốt \\ \hline
Hallucination Rate & Tỷ lệ thông tin sai so với tài liệu gốc & $<$3\% & Rất thấp \\ \hline
\end{tabularx}
\end{table}

Đặc biệt, chỉ số \textbf{Faithfulness đạt 93.8\%} là kết quả nổi bật nhất trong bộ thử nghiệm nội bộ, cho thấy chiến lược kết hợp Hybrid Search, RRF và cơ chế Reranking cấu hình được có hiệu quả trong việc cung cấp ngữ cảnh chất lượng cao cho LLM, từ đó giúp giảm hiện tượng ảo giác (Hallucination) xuống dưới 3\% trên tập kiểm thử đã xây dựng.

\rule{\textwidth}{0.5pt}

\section{So sánh hiệu quả với phương thức truyền thống}

\subsection{Phân tích tiết kiệm thời gian}

Một trong những mục tiêu cốt lõi của PoliWise là giảm thời gian tra cứu thông tin nội bộ. Nhóm thực hiện khảo sát nhỏ với 20 người dùng thử nghiệm, so sánh thời gian hoàn thành 10 nhiệm vụ tra cứu chính sách theo hai phương thức:

\begin{table}[H]
\centering
\caption{So sánh thời gian tra cứu: Truyền thống vs. PoliWise}
\label{tab:time_comparison}
\begin{tabularx}{\textwidth}{|l|X|X|X|}
\hline
\rowcolor{gray!10}
Loại nhiệm vụ & Thời gian trung bình (Truyền thống) & Thời gian trung bình (PoliWise) & Giảm (\%) \\ \hline
Tra cứu điều khoản đơn giản & 4.2 phút & 0.6 phút & 85.7\% \\ \hline
So sánh hai phiên bản chính sách & 18.5 phút & 2.1 phút & 88.6\% \\ \hline
Tổng hợp thông tin từ nhiều tài liệu & 35.0 phút & 3.8 phút & 89.1\% \\ \hline
Tra cứu quy định phòng ban cụ thể & 12.0 phút & 1.4 phút & 88.3\% \\ \hline
\textbf{Trung bình chung} & \textbf{17.4 phút} & \textbf{2.0 phút} & \textbf{88.5\%} \\ \hline
\end{tabularx}
\end{table}

Kết quả cho thấy PoliWise giúp \textbf{giảm trung bình 88.5\% thời gian tra cứu}, tương đương tiết kiệm hơn 15 phút cho mỗi nhiệm vụ. Nếu quy đổi theo báo cáo của McKinsey \cite{mckinsey2012} (9.3 giờ/tuần cho tra cứu thông tin), hệ thống có tiềm năng giải phóng hơn \textbf{8 giờ làm việc hiệu quả mỗi tuần} cho mỗi nhân viên.

\subsection{Đánh giá độ chính xác so với tra cứu thủ công}

Bên cạnh tốc độ, nhóm cũng đánh giá tỷ lệ trả lời đúng câu hỏi giữa hai phương thức:

\begin{itemize}
    \item \textbf{Tra cứu thủ công (Ctrl+F):} Người dùng trả lời đúng 74\% câu hỏi (26\% còn lại bỏ sót hoặc hiểu sai do tìm không đúng phiên bản tài liệu).
    \item \textbf{PoliWise (RAG):} Câu trả lời được xác nhận đúng và đầy đủ trong 91.2\% trường hợp (tương ứng Context Recall), đồng thời luôn kèm trích dẫn nguồn gốc để người dùng xác minh.
\end{itemize}

Sự cải thiện này đặc biệt có ý nghĩa trong bối cảnh tuân thủ pháp lý và chính sách, nơi một sai sót nhỏ trong tra cứu có thể dẫn đến rủi ro nghiêm trọng cho doanh nghiệp.

\rule{\textwidth}{0.5pt}

\section{Đánh giá phi chức năng}

\subsection{Hiệu năng hệ thống}

Kết quả kiểm thử tải nội bộ được tổng hợp và đánh giá theo ngưỡng yêu cầu. Các số liệu dưới đây dùng để phản ánh khả năng vận hành trong môi trường thử nghiệm của nhóm, chưa thay thế cho kiểm thử tải production:

\begin{table}[H]
\centering
\caption{Đánh giá hiệu năng so với yêu cầu phi chức năng}
\label{tab:nfr_evaluation}
\begin{tabularx}{\textwidth}{|l|X|X|l|}
\hline
\rowcolor{gray!10}
Tiêu chí & Yêu cầu đặt ra & Kết quả đo được & Đánh giá \\ \hline
Thời gian phản hồi & $<$ 2 giây & 1.2--1.8 giây & Đạt \\ \hline
Time-to-First-Token & $<$ 500ms & 150ms & Vượt \\ \hline
Hybrid Search latency & $<$ 100ms & 12ms & Vượt \\ \hline
Tỷ lệ uptime & 99\% & Ổn định 24/7 trong 7 ngày thử & Đạt \\ \hline
Số người dùng đồng thời & 50 người & Đến 80 người không giảm hiệu năng & Vượt \\ \hline
\end{tabularx}
\end{table}

\subsection{Đánh giá bảo mật}

Các cơ chế bảo mật của hệ thống đã được kiểm tra và xác nhận hoạt động đúng:

\begin{itemize}
    \item \textbf{Kiểm soát truy cập (RBAC):} Xác nhận nhân viên phòng ban A không thể truy cập tài liệu được giới hạn cho phòng ban B, ngay cả khi gửi câu hỏi liên quan.
    \item \textbf{Tính bảo toàn JWT:} Token hết hạn sau 15 phút, cơ chế refresh token hoạt động đúng. Khi tài khoản bị thu hồi (\texttt{Revoked}), event RabbitMQ invalidate token ngay lập tức.
    \item \textbf{Presigned URL MinIO:} Liên kết tải tài liệu gốc được cấp theo thời hạn ngắn, hiện cấu hình 1 giờ, giúp giảm rủi ro chia sẻ đường dẫn tràn lan so với URL tĩnh.
    \item \textbf{Rate Limiting:} Cơ chế giới hạn 30 câu hỏi/phút/người dùng hoạt động chính xác, trả về HTTP 429 khi vượt ngưỡng.
\end{itemize}

\rule{\textwidth}{0.5pt}

\section{Hạn chế và tồn tại}

Mặc dù đạt được nhiều kết quả khả quan, hệ thống PoliWise vẫn còn tồn tại một số hạn chế cần ghi nhận:

\begin{itemize}
    \item \textbf{Phụ thuộc chất lượng tài liệu đầu vào:} Các văn bản scan kém chất lượng hoặc bố cục phức tạp (nhiều cột, bảng biểu lồng nhau) làm giảm độ chính xác của OCR, ảnh hưởng đến chất lượng chunking và kết quả tìm kiếm.
    \item \textbf{Giới hạn ngôn ngữ:} Mô hình bge-m3 hỗ trợ đa ngôn ngữ tốt, nhưng khi văn bản trộn lẫn cả tiếng Việt, tiếng Anh và các thuật ngữ chuyên ngành, độ chính xác embedding có thể giảm nhẹ.
    \item \textbf{Yêu cầu phần cứng cao:} Để chạy đầy đủ mô hình LLM On-premise (Qwen2.5-7B), hệ thống cần GPU với 24 GB VRAM, đây là rào cản chi phí đáng kể cho các doanh nghiệp vừa và nhỏ.
    \item \textbf{Chưa hỗ trợ câu hỏi đa lượt phức tạp:} Hệ thống xử lý tốt các câu hỏi đơn lẻ, nhưng chưa triển khai cơ chế bộ nhớ hội thoại dài hạn (Long-term Conversation Memory) để xử lý các câu hỏi cần tham chiếu ngữ cảnh từ nhiều lượt trao đổi trước.
    \item \textbf{Phạm vi thực nghiệm hạn chế:} Bộ dữ liệu thử nghiệm 50 văn bản với 150 câu hỏi là chưa đủ lớn để đánh giá toàn diện khả năng mở rộng (scalability) của hệ thống trong môi trường doanh nghiệp thực tế.
\end{itemize}

\rule{\textwidth}{0.5pt}

\section{Bài học kinh nghiệm}

\subsection{Về kiến trúc và công nghệ}

\begin{itemize}
    \item \textbf{Hybrid Search là bắt buộc, không phải tùy chọn:} Trong giai đoạn đầu, nhóm chỉ triển khai Vector Search thuần túy. Kết quả thử nghiệm nội bộ cho thấy hệ thống bỏ sót các câu hỏi chứa mã số, tên viết tắt (ví dụ: ``Điều 5.2.a'' hay ``Nghị định 13/2023/NĐ-CP'') -- những trường hợp mà keyword search (BM25) vượt trội hơn hẳn. Việc bổ sung Hybrid Search + RRF \cite{rrf2009} đã cải thiện Context Recall trong tập kiểm thử của nhóm.
    \item \textbf{Reranking là lớp lọc có thể mở rộng:} Sau khi có danh sách ứng viên từ Hybrid Search, hệ thống có cơ chế gọi dịch vụ Reranker ngoài để sắp xếp lại các chunk trước khi đưa vào prompt. Trong cấu hình hiện tại, tùy chọn này được để tắt mặc định nhằm giảm độ trễ và phụ thuộc hạ tầng, nhưng vẫn được giữ như một hướng mở rộng khi cần tối ưu Faithfulness.
    \item \textbf{Phân tách dịch vụ TEI và vLLM là quyết định đúng đắn:} Khi thử nghiệm chạy chung, tiến trình embedding khối lượng lớn tài liệu cạnh tranh GPU với LLM, làm treo toàn bộ hệ thống. Việc tách thành hai container riêng biệt giải quyết hoàn toàn vấn đề này.
    \item \textbf{PostgreSQL + pgvector là lựa chọn tối ưu cho quy mô vừa:} Việc sử dụng một cơ sở dữ liệu duy nhất cho cả metadata quan hệ và vector embedding giúp giảm đáng kể độ phức tạp vận hành, đặc biệt khi thực hiện các truy vấn kết hợp RBAC (JOIN với bảng phân quyền) và tìm kiếm vector trong cùng một câu SQL.
\end{itemize}

\subsection{Về quy trình phát triển}

\begin{itemize}
    \item \textbf{Chất lượng dữ liệu quyết định tất cả:} Nhóm nhận ra rằng 80\% nỗ lực cải thiện chất lượng câu trả lời đến từ việc cải thiện pipeline tiền xử lý dữ liệu (OCR, chunking, metadata), không phải từ việc thay đổi mô hình AI.
    \item \textbf{Kiểm thử sớm với người dùng thực:} Việc đưa giao diện demo cho người dùng thật sử dụng từ sớm (Sprint 3) đã giúp nhóm phát hiện ra các vấn đề về UX mà không thể tìm thấy qua kiểm thử kỹ thuật đơn thuần, ví dụ như người dùng cần xem ngay đoạn văn bản gốc bên cạnh câu trả lời của AI.
    \item \textbf{Quản lý cấu hình tập trung:} Việc sử dụng tệp \texttt{.env} tập trung và Docker Compose giúp đơn giản hóa đáng kể quá trình chuyển đổi giữa môi trường phát triển, kiểm thử và sản xuất, đặc biệt khi cần thay đổi giữa Groq API và mô hình On-premise.
\end{itemize}

\clearpage

% =====================================================
% CHƯƠNG VI: KẾT LUẬN VÀ ĐỊNH HƯỚNG PHÁT TRIỂN
% =====================================================
\chapter{KẾT LUẬN VÀ ĐỊNH HƯỚNG PHÁT TRIỂN}

Chương này tổng kết toàn bộ quá trình thực hiện đề tài, khẳng định các đóng góp cụ thể của hệ thống PoliWise, đồng thời vạch ra lộ trình phát triển và cải tiến trong tương lai nhằm đưa hệ thống đến gần hơn với yêu cầu triển khai thực tế tại các doanh nghiệp.

\rule{\textwidth}{0.5pt}

\section{Kết luận}

\subsection{Tóm tắt những gì đã thực hiện}

Đề tài đã nghiên cứu và xây dựng thành công hệ thống Trợ lý AI hỗ trợ quản trị chính sách và tri thức nội bộ -- PoliWise -- dựa trên kiến trúc Retrieval-Augmented Generation (RAG). Cụ thể, các công việc đã hoàn thành bao gồm:

\begin{enumerate}
    \item \textbf{Nghiên cứu lý thuyết toàn diện:} Tìm hiểu sâu về các công nghệ nền tảng bao gồm kiến trúc Microservices, RAG pipeline, Vector Embedding, Semantic Search, LLM và các framework triển khai AI như LangChain, FastAPI, vLLM.

    \item \textbf{Phân tích và thiết kế hệ thống bài bản:} Khảo sát thực trạng, xác định các vấn đề cốt lõi, xây dựng đặc tả yêu cầu, thiết kế kiến trúc đa tầng với 8 microservices độc lập và thiết kế cơ sở dữ liệu gồm 14 thực thể với 16 mối quan hệ.

    \item \textbf{Triển khai hệ thống hoàn chỉnh:} Lập trình và tích hợp thành công toàn bộ các thành phần của RAG pipeline, từ OCR Fallback, Hierarchical Chunking, Embedding với BAAI/bge-m3, Hybrid Search (Vector + BM25) kết hợp RRF, cơ chế Reranking cấu hình được, đến sinh câu trả lời theo luồng (SSE Streaming) với LLM.

    \item \textbf{Kiểm thử và đánh giá:} Thực nghiệm nội bộ trên tập 50 văn bản chính sách và 150 câu hỏi kiểm thử, ghi nhận các chỉ số Ragas tham khảo: Context Precision 89.5\%, Faithfulness 93.8\%, Hallucination Rate dưới 3\% trên tập kiểm thử này.
\end{enumerate}

\subsection{Đóng góp của đề tài}

Đề tài có các đóng góp cụ thể về mặt học thuật và thực tiễn như sau:

\textbf{Về mặt học thuật:}
\begin{itemize}
    \item Trình bày một cách có hệ thống kiến trúc RAG end-to-end tối ưu cho bài toán quản trị tri thức doanh nghiệp, từ giai đoạn tiền xử lý đến giai đoạn sinh câu trả lời.
    \item Đề xuất và đánh giá chiến lược \textbf{Hierarchical Parent-Child Chunking} kết hợp với \textbf{Hybrid Search + RRF} và cơ chế \textbf{Reranking} cấu hình được như một pipeline tối ưu để cân bằng giữa độ chính xác ngữ nghĩa và khả năng tìm kiếm từ khóa chính xác.
    \item Phân tích và so sánh toàn diện các lựa chọn công nghệ (PostgreSQL+pgvector vs. Specialized Vector DB, Groq API vs. On-premise LLM) với lý luận rõ ràng dựa trên bối cảnh bài toán cụ thể.
\end{itemize}

\textbf{Về mặt thực tiễn:}
\begin{itemize}
    \item Cung cấp một mẫu tham chiếu kiến trúc (Reference Architecture) cho các doanh nghiệp vừa và nhỏ muốn triển khai hệ thống hỏi đáp AI nội bộ với chi phí thấp và bảo mật cao trên hạ tầng On-premise.
    \item Minh chứng rằng RAG với mô hình mã nguồn mở (bge-m3, LLaMA, Qwen) hoàn toàn có thể đạt chất lượng tương đương các giải pháp thương mại với chi phí vận hành thấp hơn đáng kể.
    \item Xây dựng hệ thống PoliWise với đầy đủ các tính năng cốt lõi sẵn sàng cho giai đoạn thử nghiệm pilot tại doanh nghiệp.
\end{itemize}

\rule{\textwidth}{0.5pt}

\section{Định hướng phát triển}

\subsection{Cải tiến ngắn hạn (3--6 tháng)}

Trong giai đoạn ngắn hạn, nhóm tập trung vào việc khắc phục các hạn chế đã phát hiện và tăng độ ổn định của hệ thống:

\begin{itemize}
    \item \textbf{Bộ nhớ hội thoại dài hạn (Conversation Memory):} Tích hợp cơ chế lưu trữ và nén lịch sử hội thoại (ví dụ: sử dụng kỹ thuật Summarization Memory) để hệ thống có thể xử lý các câu hỏi đa lượt phức tạp, tham chiếu ngữ cảnh từ các lượt trước.
    \item \textbf{Cải thiện OCR cho tài liệu phức tạp:} Tích hợp mô hình OCR thế hệ mới như \texttt{Surya} hoặc \texttt{Nougat} để xử lý tốt hơn các tài liệu có bố cục phức tạp (nhiều cột, bảng biểu, công thức toán học).
    \item \textbf{Tối ưu chi phí phần cứng:} Nghiên cứu áp dụng kỹ thuật lượng tử hóa mô hình (Model Quantization -- GGUF 4-bit/8-bit) để giảm yêu cầu VRAM từ 24 GB xuống 8--12 GB, mở rộng khả năng triển khai cho các GPU phổ thông hơn (RTX 3070/4070).
    \item \textbf{Tăng cường kiểm thử:} Mở rộng bộ câu hỏi kiểm thử lên 500+ câu hỏi, bổ sung các trường hợp kiểm thử tiêu cực (adversarial queries) và câu hỏi ngoài phạm vi tài liệu để đánh giá khả năng từ chối trả lời hợp lý của hệ thống.
\end{itemize}

\subsection{Phát triển trung hạn (6--18 tháng)}

Trong giai đoạn trung hạn, nhóm định hướng bổ sung các tính năng nâng cao nhằm tăng giá trị sử dụng:

\begin{itemize}
    \item \textbf{Trợ lý thông minh đa phương thức (Multimodal):} Mở rộng khả năng xử lý để có thể hiểu và trả lời câu hỏi dựa trên nội dung hình ảnh, biểu đồ và sơ đồ được nhúng trong tài liệu chính sách, sử dụng các mô hình Vision-Language như LLaVA hoặc Qwen-VL.
    \item \textbf{Tích hợp nguồn dữ liệu thời gian thực:} Xây dựng các connector kết nối trực tiếp với các hệ thống doanh nghiệp phổ biến như Microsoft SharePoint, Google Drive, Confluence và hệ thống CRM/ERP để đồng bộ tài liệu tự động, loại bỏ bước tải lên thủ công.
    \item \textbf{Hệ thống tự học (Active Learning):} Triển khai vòng lặp phản hồi tự động: dữ liệu Like/Dislike từ người dùng được dùng để tinh chỉnh (fine-tune) mô hình Reranker hoặc cải thiện chiến lược Chunking cho từng loại tài liệu đặc thù.
    \item \textbf{Hỗ trợ phân tích chính sách chuyên sâu:} Phát triển module ``Policy Analyst'' cho phép AI chủ động so sánh chính sách nội bộ với các quy định pháp luật hiện hành (ví dụ: Bộ luật Lao động, GDPR), phát hiện các mâu thuẫn hoặc khoảng trống cần bổ sung.
\end{itemize}

\subsection{Tầm nhìn dài hạn (18 tháng trở lên)}

Trong tầm nhìn dài hạn, PoliWise hướng tới trở thành một nền tảng quản trị tri thức doanh nghiệp toàn diện:

\begin{itemize}
    \item \textbf{Nền tảng SaaS đa thuê bao (Multi-tenant SaaS):} Phát triển phiên bản dịch vụ đám mây cho phép nhiều doanh nghiệp độc lập cùng sử dụng nền tảng với dữ liệu hoàn toàn tách biệt, giảm chi phí triển khai và vận hành cho từng tổ chức.
    \item \textbf{Đồ thị tri thức doanh nghiệp (Enterprise Knowledge Graph):} Xây dựng thêm tầng đồ thị tri thức (Knowledge Graph) để nắm bắt các mối quan hệ ngữ nghĩa phức tạp giữa các chính sách, quy trình và thực thể nghiệp vụ, cho phép thực hiện các truy vấn suy luận (Reasoning Query) phức tạp hơn thuần RAG.
    \item \textbf{Trợ lý AI chủ động (Proactive AI Agent):} Phát triển PoliWise từ hệ thống phản ứng (Reactive) sang chủ động (Proactive), tức là AI có khả năng tự phát hiện và cảnh báo khi phát hiện tài liệu sắp hết hạn, khi có chính sách pháp luật mới có thể ảnh hưởng đến quy định nội bộ, hoặc khi phát hiện mâu thuẫn giữa các chính sách.
    \item \textbf{Hệ sinh thái tích hợp rộng hơn:} Tích hợp với các nền tảng giao tiếp nội bộ như Microsoft Teams, Slack để nhân viên có thể tra cứu chính sách trực tiếp trong môi trường làm việc hàng ngày mà không cần mở thêm ứng dụng.
\end{itemize}

\rule{\textwidth}{0.5pt}

\section{Lời kết}

Đề tài ``PoliWise -- Hệ thống Trợ lý AI hỗ trợ quản trị chính sách và tri thức nội bộ'' đã được thực hiện với mong muốn giải quyết một bài toán thực tiễn có giá trị cao trong môi trường doanh nghiệp hiện đại: làm thế nào để nhân viên có thể tiếp cận thông tin nội bộ một cách nhanh chóng, chính xác và bảo mật trong bối cảnh khối lượng tài liệu ngày càng lớn và phức tạp.

Kết quả thực nghiệm cho thấy kiến trúc RAG kết hợp với Hybrid Search, RRF, cơ chế Reranking cấu hình được và các mô hình ngôn ngữ mã nguồn mở có khả năng đáp ứng các yêu cầu khắt khe về độ chính xác và bảo mật trong môi trường On-premise. Đặc biệt, việc tích hợp sâu phân quyền RBAC vào chính tầng tìm kiếm vector là một đóng góp kỹ thuật quan trọng, đảm bảo rằng mỗi câu trả lời của AI đều phù hợp với quyền hạn của từng người dùng cụ thể.

Nhóm hy vọng những kết quả và kinh nghiệm tích lũy trong đề tài này sẽ là nền tảng vững chắc để tiếp tục phát triển PoliWise thành một sản phẩm thực sự có ích cho cộng đồng doanh nghiệp Việt Nam, đồng thời đóng góp vào xu hướng ứng dụng AI có trách nhiệm (Responsible AI) trong lĩnh vực quản trị tri thức.

\clearpage

% =====================================================
% TÀI LIỆU THAM KHẢO
% =====================================================
\phantomsection
\addcontentsline{toc}{chapter}{TÀI LIỆU THAM KHẢO}

\renewcommand{\bibname}{TÀI LIỆU THAM KHẢO}

\begin{thebibliography}{99}

\bibitem{mckinsey2012}
M. Chui, J. Manyika, J. Bughin, R. Dobbs, C. Roxburgh, H. Sarrazin, G. Sands, and M. Westergren,
\textit{The social economy: Unlocking value and productivity through social technologies},
McKinsey Global Institute, 2012.

\bibitem{ragas2024}
S. Es, J. James, L. Espinosa-Anke, and S. Schockaert,
\textit{RAGAs: Automated Evaluation of Retrieval Augmented Generation},
In Proceedings of the 18th Conference of the European Chapter of the Association for Computational Linguistics: System Demonstrations, pages 150--158, 2024.

\bibitem{bgem32024}
J. Chen, S. Xiao, P. Zhang, K. Luo, D. Lian, and Z. Liu,
\textit{M3-Embedding: Multi-Linguality, Multi-Functionality, Multi-Granularity Text Embeddings Through Self-Knowledge Distillation},
Findings of the Association for Computational Linguistics: ACL 2024, pages 2318--2335, 2024.

\bibitem{rrf2009}
G. V. Cormack, C. L. Clarke, and S. Buettcher,
\textit{Reciprocal rank fusion outperforms threshold and joint methods for meta-search},
In Proceedings of the 32nd international ACM SIGIR conference on Research and development in information retrieval, pages 758--759, 2009.

\bibitem{fowler2002}
Martin Fowler,
\textit{Patterns of Enterprise Application Architecture},
Addison-Wesley, 2002.

\bibitem{reactdev}
React Official Documentation,
\url{https://react.dev}.

\bibitem{nestjscom}
NestJS Official Documentation,
\url{https://nestjs.com}.

\end{thebibliography}

\end{document}
